% Negative space between rows of boxes, as in real papers' figures: "\\ \vspace{-len}" between two
% rows. Under \centering the \\ ends the paragraph, and the space is a \vskip between two paragraphs,
% which pulls the second row up over the first (the interline glue is \lineskip there); in running
% text it is a \vadjust below the line.
\documentclass{article}
\newcommand\panel[1]{\parbox[b]{3cm}{\centering\rule{3cm}{1.5cm}\\ (#1) Panel}}
\begin{document}
Text before the rows, long enough to make a full line of running text in the column and a bit more.

\begin{center}
\panel{a}\hspace{-5mm}\panel{b}
\\ \vspace{-0.35cm}
\panel{c}\hspace{-5mm}\panel{d}
\end{center}
Text between, long enough to make a full line of running text in the column and a bit more.

\noindent\rule{3cm}{1cm}\hspace{-5mm}\rule{3cm}{1cm}\\ \vspace{-0.25cm}
\rule{3cm}{1cm}\hspace{-5mm}\rule{3cm}{1cm}

Running text with space taken back in the middle of a sentence, long enough to fill a line or two
of the column at every width\vspace{-2pt} it is set at, and then it goes on for another line or so.
\begin{figure}[h]
\centering
\panel{e}\hspace{-5mm}\panel{f}
\\ \vspace{-0.35cm}
\panel{g}\hspace{-5mm}\panel{h}
\caption{Two rows of panels.}
\end{figure}
Text after, long enough to make a full line of running text in the column and a bit more.
\end{document}
